%! Author = andrea
%! Date = 11/11/2025

% Preamble
\documentclass[11pt]{article}

% Packages
\usepackage{amsmath}
\usepackage{amssymb}
\usepackage{geometry}
\usepackage{graphicx}
\usepackage{hyperref}
\usepackage{enumitem}
\setlist[itemize]{noitemsep, topsep=2pt}
\setlist[enumerate]{noitemsep, topsep=2pt}
\geometry{margin=2.5cm}

% Document
\begin{document}

\begin{center}
    {\large PHY1133: Python for Physicists}\\[0.5em]
    {\large Academic Year: 2025--2026}\\[0.5em]
    {\Large Planning a Flight Route over the Mediterranean}\\[0.5em]
    {\large Dr Andrea De Marco}\\[0.5em]
\end{center}
\noindent\textbf{Submission Deadline:} [31 January 2026]\\[0.5em]

\section{Scenario}
You are helping a small cargo aircraft plan a route to deliver equipment to several airports around the Mediterranean. The pilot wants a route that
\begin{itemize}
    \item starts at a home base (Malta International Airport),
    \item visits every other airport exactly once,
    \item returns to Malta, and
    \item has the shortest possible distance along Earth's surface.
\end{itemize}

This is an instance of the Traveling Salesperson Problem (TSP) on a sphere. In this assignment you will:
\begin{itemize}
    \item work with real-looking latitude/longitude coordinates,
    \item implement the Haversine formula for distances on Earth,
    \item implement two TSP algorithms:
    \begin{enumerate}
        \item brute force search (exact, small problems),
        \item nearest neighbour heuristic (fast, approximate),
    \end{enumerate}
    \item compute and compare route lengths, and
    \item visualise the routes.
\end{itemize}

You should use Python tools you already know: variables, conditions, loops, lists, functions, basic NumPy, and simple plotting.

\section{Data: Airports Around the Mediterranean}
You are given 10 airports with approximate latitude and longitude in degrees:

\begin{verbatim}
# name, latitude (deg), longitude (deg)
airports = [
    ("Malta Int. (MLA)",     35.857,  14.477),
    ("Catania (CTA)",        37.469,  15.066),
    ("Palermo (PMO)",        38.176,  13.103),
    ("Rome Fiumicino (FCO)", 41.800,  12.238),
    ("Athens (ATH)",         37.936,  23.947),
    ("Naples (NAP)",         40.886,  14.290),
    ("Tunis (TUN)",          36.851,  10.228),
    ("Algiers (ALG)",        36.691,   3.215),
    ("Barcelona (BCN)",      41.297,   2.078),
    ("Nice (NCE)",           43.665,   7.215),
]
\end{verbatim}

Treat index \texttt{0} (\texttt{"Malta Int. (MLA)"}) as the fixed starting and ending airport. Tours will start at index 0, visit all indices 1--9 exactly once, then return to 0.

\section{Coordinates, Distances, and Physics}
\subsection{Latitude and Longitude}
\begin{itemize}
    \item Latitude (lat): north--south angle from the equator.
    \begin{itemize}
        \item 0$^\circ$ at the equator,
        \item +90$^\circ$ at the North Pole,
        \item \textminus 90$^\circ$ at the South Pole.
    \end{itemize}
    \item Longitude (lon): east--west angle from Greenwich.
    \begin{itemize}
        \item 0$^\circ$ at Greenwich,
        \item positive eastwards, negative westwards.
    \end{itemize}
\end{itemize}

Python's trigonometric functions (\texttt{math.sin}, \texttt{math.cos}, etc.) use radians, not degrees. The conversion from degrees to radians is
\[
\text{radians} = \text{degrees} \times \frac{\pi}{180}.
\]

%In Python this can be implemented as
%\begin{verbatim}
%import math
%
%def deg2rad(deg):
%    return deg * math.pi / 180.0
%\end{verbatim}

\subsection{Flat Euclidean Distance}
If we pretend the Earth is flat over this small region and use latitude and longitude as $y$ and $x$, we can define a local planar distance in \emph{kilometres} by converting degrees to km first. Let $P=(\text{lat}_1,\text{lon}_1)$ and $Q=(\text{lat}_2,\text{lon}_2)$ in degrees and define
\[
  \Delta \text{lat} = \text{lat}_2 - \text{lat}_1,\qquad
  \Delta \text{lon} = \text{lon}_2 - \text{lon}_1.
\]
Using Earth's radius $R=6371\,$km, set
\[
  k_{\text{lat}} = \frac{\pi R}{180},\qquad
  k_{\text{lon}} = k_{\text{lat}} \cos\!\bigl(\bar{\varphi}\bigr),\quad
  \bar{\varphi} = \frac{\text{lat}_1+\text{lat}_2}{2}\cdot\frac{\pi}{180} \ (\text{radians}).
\]
The local Euclidean distance in km is then
\[
  d_{\text{euclid}} = \sqrt{\bigl(k_{\text{lon}}\Delta \text{lon}\bigr)^2 + \bigl(k_{\text{lat}}\Delta \text{lat}\bigr)^2}\,.
\]
This ``flat'' approximation is not physically exact over large areas, but on this dataset it is acceptable for comparison with Haversine.

\subsection{Haversine Distance (Great-Circle Distance)}
To approximate the physical distance along Earth's surface, we use the Haversine formula.

For two points with coordinates $(\text{lat}_1, \text{lon}_1)$ and $(\text{lat}_2, \text{lon}_2)$ in degrees, we define
\begin{align*}
\Delta \text{lat} &= \text{lat}_2 - \text{lat}_1,\\
\Delta \text{lon} &= \text{lon}_2 - \text{lon}_1.
\end{align*}

After converting all angles to radians, the Haversine expression is
\[
 a = \sin^2\left(\frac{\Delta \text{lat}}{2}\right)
     + \cos(\text{lat}_1)\cos(\text{lat}_2)\sin^2\left(\frac{\Delta \text{lon}}{2}\right).
\]

The central angle $c$ is
\[
 c = 2 \cdot \arctan2\bigl(\sqrt{a}, \sqrt{1 - a}\bigr),
\]

and the Haversine distance is
\[
 d_{\text{haversine}} = R \cdot c,
\]
where $R$ is Earth's radius in km (use $R = 6371$~km). You must implement this function.

%In Python:
%\begin{verbatim}
%def haversine_distance(p, q, R=6371.0):
%    """p, q are (lat_deg, lon_deg); returns distance in km."""
%    lat1_deg, lon1_deg = p
%    lat2_deg, lon2_deg = q
%
%    lat1 = deg2rad(lat1_deg)
%    lon1 = deg2rad(lon1_deg)
%    lat2 = deg2rad(lat2_deg)
%    lon2 = deg2rad(lon2_deg)
%
%    dlat = lat2 - lat1
%    dlon = lon2 - lon1
%
%    a = (math.sin(dlat/2)**2
%         + math.cos(lat1)*math.cos(lat2)*math.sin(dlon/2)**2)
%    c = 2 * math.atan2(math.sqrt(a), math.sqrt(1 - a))
%
%    return R * c
%\end{verbatim}

\section{The Traveling Salesperson Problem (TSP)}
We have $N = 10$ airports, indexed from 0 to 9. A \emph{tour} is a list of indices
\[
 [0, i_1, i_2, \dots, i_9, 0]
\]
where 0 appears at the beginning and the end, and each of the other indices 1--9 appears exactly once. The length of a tour is the sum of distances between consecutive airports along the route, using the chosen distance metric (Haversine distance in km). The goal is to find a tour with the smallest possible total Haversine distance.

\section{Algorithms to Implement}
You must implement both of the following algorithms:
\begin{enumerate}
    \item brute force search (exact),
    \item nearest neighbour heuristic (constructive).
\end{enumerate}

\subsection{Brute Force Search}
\textbf{Strategy:}
\begin{itemize}
    \item Fix starting index 0 (Malta).
    \item Generate all permutations of the remaining indices $[1,2,\dots,9]$.
    \item For each permutation, build the route $[0] + \text{perm} + [0]$ and compute its length.
    \item Keep the shortest route and its length.
\end{itemize}

Pseudocode outline:
\begin{enumerate}
    \item $\text{start} = 0$.
    \item $\text{others} = [1,2,\dots,9]$.
    \item Set $\text{best\_route} = \text{None}$ and $\text{best\_length} = +\infty$.
    \item For each permutation \texttt{perm} of \texttt{others}:
    \begin{enumerate}
        \item Build \texttt{route = [start] + list(perm) + [start]}.
        \item Compute $L = \text{route\_length(route, airports, haversine\_distance)}$.
        \item If $L < \text{best\_length}$, update the best route and length.
    \end{enumerate}
    \item Return \texttt{best\_route} and \texttt{best\_length}.
\end{enumerate}

You may use \texttt{itertools.permutations} to generate permutations.

\subsection{Nearest Neighbour Heuristic}
\textbf{Strategy:}
\begin{itemize}
    \item Start at Malta.
    \item At each step, move to the closest unvisited airport.
    \item After visiting all airports, return to Malta.
\end{itemize}

Pseudocode:
\begin{enumerate}
    \item $\text{start} = 0$.
    \item $\text{current} = \text{start}$.
    \item $\text{route} = [\text{current}]$.
    \item $\text{unvisited} = \{1,2,\dots,9\}$.
    \item While \texttt{unvisited} is not empty:
    \begin{enumerate}
        \item For each candidate in \texttt{unvisited}, compute the distance from \texttt{current}.
        \item Choose the candidate with the smallest distance.
        \item Append it to \texttt{route}, remove it from \texttt{unvisited}, and set \texttt{current = candidate}.
    \end{enumerate}
    \item Append \texttt{start} to \texttt{route} to close the loop.
    \item Return \texttt{route}.
\end{enumerate}

You should then compute the length of this route using Haversine distance.

\section{Programming Tasks}
Submit a single \texttt{.py} file or a Jupyter notebook.

\subsection*{Task 1: Data and Helpers}
\begin{enumerate}
    \item Store the \texttt{airports} list as given.
    \item Implement a helper to get coordinates by index.
\end{enumerate}

\subsection*{Task 2: Distance Functions}
\begin{enumerate}
    \item Implement \texttt{deg2rad(deg)}.
    \item Implement
    \begin{verbatim}
    def euclidean_distance(p, q, R=6371.0):
        ...
    \end{verbatim}
    where \texttt{p} and \texttt{q} are \texttt{(lat\_deg, lon\_deg)}; return the local flat-plane distance \emph{in kilometres} using the small-area conversion
    $k_{\text{lat}}=\pi R/180$ and $k_{\text{lon}}=k_{\text{lat}}\cos(\bar{\varphi})$ with $\bar{\varphi}$ the mean latitude in radians (see the ``Flat Euclidean Distance'' subsection).
    \item Implement
    \begin{verbatim}
    def haversine_distance(p, q, R=6371.0):
        ...
    \end{verbatim}
    where \texttt{p} and \texttt{q} are \texttt{(lat\_deg, lon\_deg)}.
    \item Implement a route length function
    \begin{verbatim}
    def route_length(route, airports, distance_fn=haversine_distance):
        """route is a list like [0, 3, 5, ..., 0]; returns total distance in km."""
    \end{verbatim}
    looping over consecutive index pairs in \texttt{route}. Use the \texttt{distance\_fn} argument so you can pass either \texttt{haversine\_distance} or \texttt{euclidean\_distance}.
    \item Optional: compare route lengths using both metrics (Euclidean vs Haversine) and comment briefly.
\end{enumerate}

\subsection*{Task 3: TSP Algorithms}
\begin{enumerate}
    \item Implement brute force:
    \begin{verbatim}
    def brute_force_route(airports, start_index=0):
        """Returns (best_route, best_length) using exhaustive search."""
    \end{verbatim}
    \item Implement nearest neighbour:
    \begin{verbatim}
    def nearest_neighbour_route(airports, start_index=0):
        """Returns a route using the nearest neighbour heuristic."""
    \end{verbatim}
    \item Run both algorithms:
    \begin{itemize}
        \item brute force to get an exact best route \texttt{(route\_exact, L\_exact)},
        \item nearest neighbour to get a heuristic route \texttt{(route\_nn, L\_nn)}.
    \end{itemize}
    Print and compare both route lengths.
\end{enumerate}

\section{Visualisation}
Using \texttt{matplotlib}:
\begin{enumerate}
    \item Plot all airports as points with
    \begin{itemize}
        \item longitude on the $x$-axis,
        \item latitude on the $y$-axis,
        \item labels for each point with its name.
    \end{itemize}

\item Plot two routes:
\begin{itemize}
    \item the exact route (brute force),
    \item the nearest neighbour route.
\end{itemize}
    You may either use different line styles/colours on the same figure or separate figures.
\end{enumerate}

\section{Short Written Answers}
Answer briefly (a few sentences each) in your write-up:
\begin{enumerate}
    \item In your own words, what is the Traveling Salesperson Problem?
    \item Why is the Haversine formula more appropriate than a simple flat Euclidean distance for these coordinates?
    \item Compare the two algorithms (brute force and nearest neighbour) in terms of
    \begin{itemize}
        \item quality of the solution on this dataset (which route is shortest), and
        \item what you expect to happen as the number of airports becomes very large.
    \end{itemize}
    \item If you experimented with another distance function, did the ``best'' route change?
\end{enumerate}

\section{Marking Scheme and Deliverables}
The assignment also includes a write-up, which should be about 3 to 4 pages long, including figures and plots, but excluding any code listing. The write-up should include:
\begin{itemize}
    \item A discussion of the problem and how you tackled it.
    \item An explanation of how the various aspects of your code works.
    \item Plots of the various route solutions.
\end{itemize}
Submit your source code, as a Python file or as a Colab Notebook along with your report in a ZIP file on the VLE.

The marking breakdown is as follows:
\begin{itemize}
    \item Data representation and appropriate functions with separation of concerns: 15\%.
    \item Correct distance function implementation and route length: 20\%.
    \item Brute force implementation and result: 15\%.
    \item Nearest neighbour implementation and result: 15\%.
    \item Visualisation and clarity of plots: 15\%.
    \item Written report with answers and interpretation: 20\%.
\end{itemize}

\end{document}